\documentclass[AER]{AEA}
\usepackage{graphicx} %package to manage images
\graphicspath{ {./figures/} }
\usepackage{hyperref}
\usepackage{caption}
\usepackage[font=scriptsize]{subcaption}
\captionsetup[figure]{labelsep=none}
\captionsetup[table]{labelsep=none}
\usepackage{bbm}
\usepackage{amsmath}
\usepackage{import}
\usepackage{array}
\usepackage{booktabs}
\usepackage{afterpage}
\usepackage{floatrow}
\usepackage{pdflscape}
\usepackage{soul}
\usepackage{float}
\usepackage{adjustbox}
\usepackage{longtable}
\usepackage{caption}
\usepackage{setspace}
\usepackage{afterpage}





\draftSpacing{1.5}

\begin{document}
\title{}
\shortTitle{}
\author{April 28, 2025}
\date{\today}

\begin{abstract}
Maternal undernutrition is an important cause of child mortality and morbidity. Prior research suggests that maternal undernutrition, as measured by pre-pregnancy underweight, is common in India. This paper is the first to document differences in maternal undernutrition by Indian social group. We apply reweighting decomposition methods to the National Family Health Survey (2019-21) data to document that pre-pregnancy underweight is about 70\% higher among Adivasis and 40\% higher among Dalits than it is among forward caste Hindus and Muslims. Although high fertility and low birth spacing are known to contribute to maternal undernutrition in other contexts, we show that these factors do not explain the gaps that we document. Instead, accounting for differences in economic status explains a large fraction of the pre-pregnancy underweight gap between Dalits and Forward Castes and between Adivasis and Forward Castes. Accounting for wealth differences reveals that despite multidimensional disadvantage, Muslims have lower rates of pre-pregnancy underweight than forward caste Hindus.

\end{abstract}

\maketitle

\section{Introduction}

Maternal undernutrition is important for population health. The period before and during pregnancy are especially important points for intervention. [EXPAND]

India is a unique case with disproportionately poor health outcomes for women of reproductive age [CITE]. This can be linked to its disproportionately poor child health outcomes [CITE].

Data on the state of maternal health, specifically, is more sparse. Estimates of pre-pregnancy health as opposed to using WRA as a proxy are important [ASK DIANE].

Using a novel reweighting method, Coffey (2015) estimates that 42.2\% of women in India entered pregnancy underweight, and Indian women gain only about 7 kg during pregnancy.

This paper produces more recent estimates using the NFHS-5 and also the first estimates of pre-pregnancy underweight by social group.

Disaggregation by social group is important because disparities in maternal health and nutrition across caste and religion are well documented. National averages mask stark inequalities [CITE].

We find that 26.5\% of Adivasi women and 22.3\% of Dalit women were underweight before pregnancy, compared to 15.5\% of forward-caste Hindus and Muslims. This represents a 70\% higher rate among Adivasis and a 43\% higher rate among Dalits. 

This paper also contributes to a growing literature on the determinants of social group disparities in health. [CITE] [READ] [ASK DIANE]

The disparities in pre-pregnancy underweight we observe are not explained by wealth, education, or fertility patterns alone: Kitagawa decompositions show that over 85\% of the Adivasi–forward and Dalit–forward caste gaps remain even after accounting for differences in parity, birth spacing, and wealth quartile.

Accounting for wealth differences reveals that despite multidimensional disadvantage, Muslims have lower rates of pre-pregnancy underweight than forward caste Hindus. [READ]





\bibliographystyle{apalike}
\bibliography{writing/references}

\end{document}
